% Chapter 9 - Conclusion
\section{Conclusion}
\label{sec:conclusion}

\subsection{Progress and Contributions}
\label{sec:progress_contributions}
This thesis presents the design, implementation, and validation of an automated pipeline for generating, validating, and evaluating road network maps from OpenStreetMap (OSM) data in the CARLA simulator.
The primary contribution is not only automated map generation, but a reproducible and auditable experimental framework for quantifying structural domain gap and validating perception-capture capability under controlled conditions.

\paragraph{Pipeline reliability and auditability.}
A complete OSM$\rightarrow$OpenDRIVE$\rightarrow$CARLA pipeline was implemented to support large-scale map generation, tiling, validation, and downstream perception evaluation.
The pipeline was hardened to ensure experimental integrity: failures propagate reliably (non-zero exit codes), misleading ``successful completion'' states were eliminated, and each run produces explicit diagnostic artifacts.
This prevents silent failures from contaminating quantitative results.

\paragraph{Tile-based validation and diagnostics.}
Fine-grained tile-level validation was introduced to localize structural defects and enable scalable quality assessment.
Reporting was redesigned to provide unambiguous counts of passed, failed, and unknown tiles, alongside a failure taxonomy that reveals dominant error modes (e.g., import failures, spawn failures, topology errors).
These diagnostics support statistically meaningful comparisons across generated maps and across pipeline configurations.

\paragraph{Determinism and reproducibility methodology.}
The thesis explicitly separates variability from (i) OSM acquisition, (ii) OSM-to-OpenDRIVE conversion, and (iii) OpenDRIVE-to-CARLA import/normalization.
The pipeline supports hashing and semantic signatures to distinguish formatting differences from true structural nondeterminism, enabling determinism audits beyond naive file equality checks.

\paragraph{Manual baseline integration.}
A manually authored CARLA map of Ingolstadt was extracted in OpenDRIVE format as a structural ground-truth baseline.
This enables direct structural comparison against automatically generated maps and supports quantitative domain gap analysis of lane topology, connectivity, and geometric properties.

\paragraph{Sensor calibration correctness for perception evaluation.}
Critical extrinsic-handling errors were identified and corrected to ensure mathematically consistent camera and LiDAR placement in the simulator.
A verification procedure is integrated to prevent systematic calibration bias, which is essential for perceptual comparisons and downstream model evaluation.

\paragraph{Summary.}
In summary, this thesis contributes (i) a reproducible map generation and evaluation pipeline, (ii) tile-based structural validation with actionable diagnostics, (iii) a determinism auditing methodology grounded in hashes and semantic signatures, (iv) integration of manual maps as structural baselines, and (v) a corrected, verifiable sensor calibration pipeline for perception analysis.

\paragraph{Primary structural domain-gap findings.}
The authoritative structural comparison (run\_11, governed) between the auto-generated and manually authored Ingolstadt maps yields:
whole-network geometry RMSE $= 54.84\,$m;
matched-subset RMSE $= 24.47\,$m for the 73.0\% of sampled road points within a 50\,m criterion;
Hausdorff distance $= 1{,}663\,$m.
These three values must be read together: the whole-network RMSE includes an unmatched-road penalty from the 4.88$\times$ road-length surplus of the auto map; the matched-subset RMSE isolates local alignment quality for the portion of the manual network that has a close auto counterpart. The road-length ratio is taken from the audited governed run\_11 coverage context, while junction/connectivity summaries come from the governed connectivity companion artifact family.
The corrected curvature support measurement on the same authoritative manual-vs-contract-run pair shows that the curvature distribution of the auto map is substantially wider than the manual reference (std $= 0.865$ vs.\ $0.055\,\mathrm{m}^{-1}$), indicating high-curvature outlier segments introduced by procedural paramPoly3 generation.
Lane and junction connectivity analysis shows that the auto map carries no road-level predecessor/successor link declarations (relying entirely on junction-based connectivity), while the manual map has 93--97\% road-link coverage at 100\% validity.
Junction laneLink completeness is 100\% in both maps.
All geometry metrics are planar (2D) because the final analysis artifact remained flat; elevation-gap quantification is deferred to future work.

\subsection{Additional contributions enabling thesis-grade evidence}
\label{sec:additional_contributions}
Beyond the core pipeline described above, the following additions strengthen experimental validity and widen metric coverage:

\begin{itemize}
  \item \textbf{Deterministic, traceable execution.} Each run records a settings snapshot and cryptographic hashes for key inputs and outputs, enabling a per-run reproducibility fingerprint that supports auditability across machines and reruns.
  \item \textbf{Explicit refusal logic for invalid tile analysis.} When tile correspondence is empty or coordinate frames are incompatible, the pipeline refuses per-tile reporting rather than producing misleading numbers. This transforms a potential failure mode into an explicit validity guarantee.
  \item \textbf{Determinism tooling.} Dedicated comparison scripts evaluate repeated runs under identical inputs to distinguish semantic nondeterminism from superficial variability (e.g., ordering or timestamps).
  \item \textbf{Multi-category OpenDRIVE-only structural metrics.} In addition to geometric distance measures (RMSE, Hausdorff), the evaluation includes: curvature distribution gap (paramPoly3-aware sampling, KL divergence, and standard deviation comparison); semantic object gap (object-type counts, road-type lengths, traffic-light and lane-marking density); topology gap (road and junction counts); and a lane/junction connectivity gap suite covering road-link predecessor/successor validity rates, lane-count distributions, junction degree, and lane-link completeness. These metrics are XML-only, CARLA-independent, and deterministic.
\end{itemize}

Collectively, these contributions enable trustworthy experimental analysis of map-induced domain gaps in autonomous driving simulation.

\section{Future Work}
\label{sec:future_work}
The core contribution of this thesis is the characterization of map-induced domain gaps and their reproducibility, rather than achieving state-of-the-art perception accuracy.
Several natural extensions can strengthen the empirical scope without changing the methodological foundation:

\begin{itemize}
  \item \textbf{Large-scale training and evaluation on HPC:} training across many auto-generated maps, larger datasets, and cross-city experiments (sim$\rightarrow$sim at scale).
  \item \textbf{Broader environmental sweeps:} systematic weather, lighting, and traffic-condition variation to disentangle map effects from appearance effects.
  \item \textbf{Expanded sim$\rightarrow$real evaluation:} incorporation of labeled real-world datasets or semi-supervised approaches to move beyond proxy metrics.
  \item \textbf{Richer structural alignment:} improved junction matching and alignment-based structural metrics to better explain localized perceptual shifts.
\end{itemize}

\paragraph{Repo-backed extension roadmap.}
This thesis quantifies the map-domain gap but does not yet close the loop with active mitigation. Existing governed scaffolding already points to several realistic extensions: reinforcement-learning-style map fuzzing for controlled perturbation of lane geometry and object density, LLM-assisted OpenDRIVE semantic checking, broader scenario and route sweeps, and realism-oriented augmentation modules. These components were not part of the evaluated thesis evidence bundle, but they define a concrete path from passive gap measurement toward robustness-oriented stress testing and automated map validation.

\paragraph{Elevation domain gap quantification.}
A later governed supplementary run demonstrated end-to-end elevation capability by producing a non-flat DEM-backed Ingolstadt XODR, but it was not adopted as the primary thesis artifact because vertical continuity checks still reported anomalies (805/7966 links above 0.5\,m; max $\Delta z \approx 6.53$\,m) and the thesis evidence bundle had already frozen around the planar run\_11 authority. The final structural comparison reported in this thesis therefore remains planar-only.

\paragraph{RQ3--RQ5: perception-driven evaluation.}
Paired perception data from Grid0828 and the auto-generated Ingolstadt map remains the prerequisite for perceptual-gap quantification and any downstream transfer validation. The collection and training entrypoints are implemented, but the thesis does not claim those experiments were completed.

\paragraph{NPC traffic density as a perception gap dimension.}
The \texttt{--num-npcs} control is wired for generated-tile collection, but the effect of dynamic traffic density on occlusion and sensor failure modes is deferred to dedicated ablation experiments.

\paragraph{Grid0828 CARLA compatibility.}
Grid0828-specific map-name normalization and cooked-world loading stability remain an open runtime compatibility issue. Current thesis-facing perception reference runs therefore remain bounded and incomplete.

Importantly, these extensions are optional: the thesis results remain valid under the controlled experimental protocol presented here, and future work primarily increases statistical power and coverage.

\section{Limitations and Failure Modes}

\subsection{Threats to validity and mitigations}
\label{sec:threats_mitigations}
Several failure modes can silently invalidate quantitative results if not handled explicitly. This thesis therefore treats the following as first-class threats, with documented mitigations:

\begin{itemize}
  \item \textbf{Alignment fallback can corrupt geometric metrics.} If an alignment stage fails and downstream metrics fall back to an identity transform, geometric distances (e.g., RMSE or Hausdorff proxies) can become meaningless. The evaluation includes a strict ``kill switch'' that disables metric reporting unless alignment status is explicitly confirmed (or an override is intentionally enabled and recorded).
  \item \textbf{Tile mismatch or empty correspondence invalidates per-tile metrics.} When correspondence is empty (header-only \texttt{correspondence.csv}) or frames are incompatible, per-tile reporting is refused by default. Any override is recorded in the run metadata.
  \item \textbf{Randomized spawning reduces repeatability.} Stochastic simulator elements (e.g., spawn placement) are controlled via explicit seeding and are logged in run metadata to ensure repeatable experiment execution.
  \item \textbf{Simulator import and normalization can alter structure.} CARLA may drop or normalize elements during import. As a fidelity check, the methodology compares coarse road, junction, and lane statistics before and after import; extending this into deeper structural equivalence checks is treated as future work.
\end{itemize}

\paragraph{Spatial correspondence validity.}
Incorrect tile pairing due to grid-origin or coordinate-frame mismatch can invalidate per-tile metrics.
To mitigate this, tile correspondence is established via spatial matching with IoU gating and optional robust translation estimation (Section~\ref{sec:tile_correspondence}).
Per-tile comparisons are refused when incompatible coordinate units are detected, and the refusal reason is recorded in the run artifacts.

\paragraph{Implication for reported results.}
When manual-to-auto tile correspondence cannot be established (e.g., zero valid matches after spatial matching), per-tile metrics are not reported.
In such cases, whole-map structural metrics remain valid and are presented as the primary evidence for structural domain gap.
The four structural geometry results must be read as a set:
whole-network RMSE $= 54.84\,$m, which includes an unmatched-road penalty arising from the 4.88$\times$ road-length surplus of the auto network;
matched-subset RMSE $= 24.47\,$m, which reflects local alignment quality for the 73.0\% of manual road samples that have an auto road counterpart within a 50\,m criterion;
and Hausdorff distance $= 1{,}663\,$m, which captures the worst-case point deviation at whole-network scale.
The whole-network RMSE and Hausdorff values are not indicators of local registration accuracy; they are dominated by roads present in one map but absent in the other.
The matched-subset RMSE is the appropriate metric for judging local geometric correspondence in the region of spatial overlap.

\label{sec:limits_failures}
This thesis emphasizes experimental control and reproducibility, yet several limitations remain and are important for interpreting results.

\subsection{Limitations}
\begin{itemize}
  \item \textbf{Perceptual Gap Comparison (RQ3):} The authoritative generated-map strict-evidence attempt targeted 50 frames, but the final governed evidence pack remained incomplete and closed with zero recorded frames. Earlier 241-frame development-run references are not used as primary governed thesis evidence. A separate Town10HD\_Opt stage probe confirmed rig callback readiness on the built-in map (\texttt{ok=true}), but that is a proxy baseline rather than paired Ingolstadt evidence. Grid0828 manual reference collection is still pending, so perceptual gap quantification requires future paired captures. A governed runtime-QA support path was fully implemented and all identified repo-local defects were resolved (wrong-map masking, bypassed streaming readiness gate, one-success false-positive stream check, None-client at \texttt{apply\_batch\_sync}, and current-world no-probe routing). In the final governed attempt on the already-loaded cooked Grid0828 world (identity confirmed, 11\,634 spawn points), the rig-attach path was entered with a valid client; however, \texttt{spawn\_response\_count}\,=\,0 with persistent streaming transport refusal yielded zero frames. This is classified as a CARLA runtime boundary during live sensor spawn on the custom cooked map rather than a remaining code defect, and is treated as an engineering boundary rather than thesis evidence.
  \item \textbf{Generalization (RQ4/RQ5 - GNN latent gap):} An earlier positive-only cosine-loss interpretation (cosine similarity 0.9989) was later rejected as a collapsed encoder. The governed contrastive artifacts instead indicate a non-collapsed configuration, with cross-cosine $=-0.160$, training cosine similarity $=-0.206$, final loss $=2.477$, and an overall verdict of \texttt{NOT\_COLLAPSED}. These results are observational only, and no downstream transfer experiment was executed, so no transfer-feasibility claim is made \cite{chen2020simclr}.
  \item \textbf{Elevation Data:} Primary thesis metrics remain planar (2D) and are anchored to run\_11. A later supplementary governed 3D capability run demonstrated a non-flat DEM-backed final XODR, but it does not replace run\_11 because vertical continuity anomalies remained and the thesis evidence bundle was already frozen.
  \item \textbf{Simulator dependence:} results are tied to CARLA 0.9.16; changes in simulator versions may alter map import behavior and sensor implementations.
\end{itemize}

\subsection{Failure modes}
Operational failure modes encountered during development are treated as part of the scientific evidence chain because they can silently invalidate results if not controlled.
Key failure modes and mitigation controls include content or version mismatches, disk-space constraints, hidden log locations, and rendering pressure on large tiled maps.

\subsection{What this method cannot guarantee}
Even with pinned inputs and explicit manifests, some sources of variability cannot be fully eliminated (e.g., low-level floating-point behavior or third-party library nondeterminism).
Therefore, conclusions are stated in terms of measurable distributions and bounded interpretation rather than assuming perfect determinism in all conditions.
